\section{Rank and payoff for every arrival order}\label{sec:accounting}

We now pass from the nominal rank guarantee of Lemma~\ref{lem:transfer}
to the actual greedy reward. A rank statistic on the final layers will
give a lower bound for every arrival order with all coins fixed. We then
compare final and nominal layers using one generator budget over the
entire auxiliary timeline. Only after these deterministic comparisons do
we average over the independent test thinning and the other masks.

\subsection{A simultaneous rank lower bound}

Fix all weights and coins. For a listed group $h$, write
$S_{h,b}=P_{h,b}\setminus O$, and define the order-independent statistic
\begin{equation}\label{eq:actual-rank-statistic}
 \lambda_h=\sum_b\rk\bigl(T_h\cap S_{h,b}\mid
                   \widehat F(b-2)+K_{h,b}\bigr).
\end{equation}
Sums over $b$ use the chosen parity; terms outside the finite nontrivial
range are zero. Each summand measures the rank left among tested,
unobserved labels of group $h$ after contracting the preceding final flat
and all lower-weight test competitors. It depends on the coins but not
on the arrival order. Let $N_{\ge h}(\pi)$ be the number of accepted labels whose
rounded weight is at least that of $G_h$.

\begin{lemma}[Order-uniform rank bound]\label{lem:all-orders}
For each realization of all coins and every permutation $\pi$,
\begin{equation}\label{eq:all-orders}
 N_{\ge h}(\pi)\ge\lambda_h.
\end{equation}
\end{lemma}
\begin{proof}
At endpoint $b$, let $F=\widehat F(b-2)$ and let $E_b$ be the entire
eligible set of that layer. Whatever the arrival order or the interleaving
with other layers, its final greedy set $B_b$ is a basis of $E_b$ over $F$.
Partition $B_b=L_b\sqcup H_b$ into selected labels of strictly lower
weight than $G_h$ and the remaining labels. Independence and spanning give
\[
 |H_b|=\rk(H_b\mid F+L_b)=\rk(E_b\mid F+L_b).
\]
Every element of $T_h\cap S_{h,b}$ outside $F$ is eligible. Those in $F$
contribute no rank. Also $L_b\subseteq K_{h,b}$, because each lower-weight
accepted label is test-positive and has nominal birth $b$ in a lower
listed group. Thus
\[
 |H_b|\ge\rk(T_h\cap S_{h,b}\mid F+L_b)
        \ge\rk(T_h\cap S_{h,b}\mid F+K_{h,b}).
\]
Sum over the layers to prove \eqref{eq:all-orders}. Including unselected
or sacrificed labels in $K_{h,b}$ only weakens this bound.
\end{proof}

\subsection{One budget for later path expansions}

\begin{lemma}[Telescoping rank cost]\label{lem:rank-cost}
For each listed group $h$, the nominal statistic $z_h$ satisfies
\begin{equation}\label{eq:rank-cost}
 \sum_b\rk(S_{h,b}\mid\widehat F(b-2)+K_{h,b})\ge z_h-c_h,
 \qquad
 c_h=|C_h|+\sum_{j>h}\bigl(|C_j|+|D_j|/\kappa\bigr).
\end{equation}
\end{lemma}
\begin{proof}
Choose the generator sets of Lemma~\ref{lem:generators} and put
\[
 J=C_h\cup\bigcup_{j>h}(C_j\cup J_j).
\]
This single set satisfies $|J|\le c_h$ and, for every integer $k$,
\begin{equation}\label{eq:common-generator-set}
 X_h(k)\subseteq\widehat F(k)\subseteq\cl(X_h(k)+J).
\end{equation}
Indeed the first group's guards are contained in $C_h$; at each later
group, its whole density path uses portions of $J_j$, and its guards
are contained in $C_j$. Induction over groups proves the second inclusion
for all times at once. The set $J$ may depend on the realized masks.

At endpoint $b$, abbreviate $A=X_h(b-2)+K_{h,b}$ and $S'=S_{h,b}$.
The conditioning set in the left side of \eqref{eq:rank-cost} lies
between $A$ and $\cl(A+J)$. The resulting rank loss is at most
\begin{align}
 \rk(S'\mid A)-\rk(S'\mid A+J)
 &=\rk(J\mid A)-\rk(J\mid A+S')\notag\\
 &\le\rk(J\mid X_h(b-2))-\rk(J\mid X_h(b)).
 \label{eq:rank-telescope-step}
\end{align}
The inequality follows from two applications of decreasing conditional
rank along the nested sets
\[
 X_h(b-2)\subseteq A\subseteq A+S'\subseteq X_h(b):
\]
the last inclusion follows from $K_{h,b}\subseteq F_{h-1}(b)$ and
$P_{h,b}\subseteq X_h(b)$. Specifically,
$\rk(J\mid A)\le\rk(J\mid X_h(b-2))$ and
$\rk(J\mid A+S')\ge\rk(J\mid X_h(b))$.
Summing the right side of
\eqref{eq:rank-telescope-step} over the chosen parity telescopes to at most
$\rk(J)\le |J|$. Missing endpoints may be added, since all differences
are nonnegative. This proves \eqref{eq:rank-cost}; neither independence
of $J$ nor a separate budget for each layer is required.
\end{proof}

\subsection{Independent thinning and group estimates}

The two preceding lemmas apply simultaneously to every realization and
arrival order. We now average the focal group's test bits, which played
no role in constructing its paths or lower-weight competition sets.

\begin{lemma}[Independent test thinning]\label{lem:thinning}
For a listed group $h$, let $\mathcal F_{-h}$ be generated by $H,D,C$,
the parity $\varepsilon$, and all test-mask memberships outside $G_h$.
Then
\begin{equation}\label{eq:test-thinning}
 \Ex[\lambda_h\mid\mathcal F_{-h}]
 \ge t\sum_b\rk(S_{h,b}\mid\widehat F(b-2)+K_{h,b})
 \ge t(z_h-c_h).
\end{equation}
\end{lemma}
\begin{proof}
The group $G_h$ is determined by $H,D$ and the fixed weights. All paths
depend on $H,D,C$ and not on $T$. The birth sets $P_{h,b}$ also use the
parity; $O=D\cup C$ omits $T$, and $K_{h,b}$ uses only lower groups' test
bits. Thus these sets, as well as $z_h$ and $c_h$, are
$\mathcal F_{-h}$-measurable, while the bits of $T_h$ remain independent
Bernoulli bits of rate $t$.

Conditional on $\mathcal F_{-h}$, choose a fixed independent witness for
each unthinned rank in \eqref{eq:test-thinning}. Its intersection with
$T_h$ is still independent over the same base and has expected size $t$
times its size. Linearity proves the first inequality, and
Lemma~\ref{lem:rank-cost} proves the second. No selected order is
conditioned upon.
\end{proof}

Return to true weight indices. For a listed true group $U_i=G_h$, let
$\lambda_i=\lambda_h$ and $c_i=c_h$; if it is unlisted, set both to zero,
as with $z_i$. Conditional on $H$, dominate the cost pointwise by counts
on all true groups, listed or not:
\[
 c_i\le |C\cap U_i|+
      \sum_{\ell>i}\left(|C\cap U_\ell|+\frac{|D\cap U_\ell|}{\kappa}\right).
\]
Only then take expectations. The mask probabilities are unchanged under
this conditioning, so
\begin{equation}\label{eq:expected-cost}
 \Ex[c_i\mid H]\le t n_i+(t+1/\kappa)\sum_{\ell>i}n_\ell.
\end{equation}
The exact density-mask contribution would be $1/(4\kappa)$; the larger
coefficient is convenient. Lemmas~\ref{lem:transfer} and
\ref{lem:thinning} now imply, for every true group of size
$n_i\ge\kappa$,
\begin{equation}\label{eq:group-payoff}
 \Ex[\lambda_i\mid H]\ge\frac t\kappa
 \left(\gamma|Y_i|-(2^{-23}+t\kappa)n_i
       -(1+t\kappa)\sum_{\ell>i}n_\ell\right).
\end{equation}
Although the different $\lambda_i$ may depend on one another's test bits,
each separate thinning calculation is valid. Their expectations can be
summed without asserting independence.

\subsection{The weighted sum and the final constant}

Let $A_{\mathrm{main}}(R,\pi)$ be the main rule's accepted set and let
$\mathcal A(H)=\{i:n_i\ge\kappa\}$ be the analyzed true groups. For a true
weight index $i$, let $N_i(\pi)$ count accepts of weight at least $B^i$.
Lemma~\ref{lem:all-orders}, with $\lambda_i=0$ on unlisted groups, gives
$\lambda_i\le N_i(\pi)$ simultaneously for every $\pi$. Each accepted
weight $B^\ell$ contributes at most
$\sum_{i\le\ell}B^i=B^\ell B/(B-1)<2B^\ell$ to the cumulative weighted
counts. Hence the pointwise inequality is
\begin{equation}\label{eq:minimum-before-expectation}
 \min_\pi u(A_{\mathrm{main}}(R,\pi))
       \ge\frac12\sum_{i\in\mathcal A(H)} B^i\lambda_i.
\end{equation}
The order has been minimized before any expectation is taken.

Multiply \eqref{eq:group-payoff} by $B^i$ and sum over analyzed groups.
The higher-weight cost is bounded using
\[
 \sum_{i\in\mathcal A(H)}B^i\sum_{\ell>i}n_\ell
 \le\sum_\ell n_\ell\sum_{i<\ell}B^i
 =\frac1{B-1}\sum_\ell B^\ell n_\ell.
\]
With the constants in \eqref{eq:constants},
\begin{equation}\label{eq:error-constant}
 2^{-23}+t\kappa+\frac{1+t\kappa}{B-1}<2^{-21}.
\end{equation}
For example $t\kappa=2^{-40}$ and the fractional term is less than
$2^{-30}$; these bounds already give the displayed strict inequality.
Combining with \eqref{eq:minimum-before-expectation} yields
\begin{equation}\label{eq:conditional-main-payoff}
 \Ex\left[\min_\pi u(A_{\mathrm{main}}(R,\pi))\given H\right]
 \ge\frac{t}{2\kappa}\left(
       \gamma\sum_{i\in\mathcal A(H)}B^i|Y_i|-2^{-21}u(U)\right).
\end{equation}

Suppose $W>0$, so $m_*>0$. Ignoring the smaller true groups loses at most
\begin{equation}\label{eq:small-groups}
 \sum_{i\notin\mathcal A(H)}B^i|Y_i|
 \le\kappa\sum_{B^i\le m_*}B^i<2\kappa m_*.
\end{equation}
If $m_*\le W/(8\kappa)$, Lemma~\ref{lem:filter} and
\eqref{eq:small-groups} give
\[
 \Ex_H\sum_{i\in\mathcal A(H)}B^i|Y_i|\ge W/4,
 \qquad \Ex_H u(U)\le W.
\]
Since $\gamma=2^{-16}$, we have
$\gamma/4-2^{-21}=7\cdot2^{-21}\ge\gamma/8$. Therefore
\begin{equation}\label{eq:main-branch-guarantee}
 \Ex\big[\min_\pi u(A_{\mathrm{main}}(R,\pi))\big]
       \ge \frac{t\gamma}{16\kappa}W.
\end{equation}
If instead $m_*>W/(8\kappa)$, the maximum branch gives more than
$W/(32\kappa)$ by \eqref{eq:maximum}.

\begin{proof}[Proof of Proposition~\ref{prop:hidden} and Theorem~\ref{thm:main}]
Choose the two branches with equal probability. In the first case above,
the main branch contributes at least $t\gamma W/(32\kappa)$ after mixing.
In the second case, the maximum branch contributes more than
$W/(64\kappa)$, which is larger than the same target. The minimum over
orders is taken separately for every complete seed, including the branch;
the fair mixture is therefore valid against branch-dependent orders.
True accepted weights dominate their rounded weights. By
\eqref{eq:rounding}, the resulting fixed-vector reward fraction is
\[
 \frac{t\gamma}{32\kappa B}
      =2^{-140-16-5-100-32}=2^{-293}.
\]
When $W=0$, the benchmark is zero and the conclusion is immediate.
Feasibility and measurability follow from Lemma~\ref{lem:feasibility} and
the analogous immediate properties of the maximum branch. This proves
Proposition~\ref{prop:hidden}. The revealed mask is selected from the
initial seed, independently of the weights, so Lemma~\ref{lem:simulation}
proves Theorem~\ref{thm:main}, with slack in its announced constant.

The construction covers empty ground sets, loops, rank zero, ties, and
zero values as specified above. Each realized weight vector is finite,
even when the distributions have unbounded support. Only the expected
optimum needs to be finite; loop values, which can never be accepted,
need no separate moment condition.
\end{proof}
